\chapter{从十四分卷到单文件源码与证明过程}
\label{chap:source-delivery}
\section{本轮使用的完整源码版本}
作者公告脚注列出 \path{lean.tar.gz.part01} 至 \path{lean.tar.gz.part14}。它们是同一个 gzip 归档的分卷，需要先按编号拼接，再解压；每个分卷并不是一个独立 Lean 模块。

本轮原站首分卷请求返回 HTTP 403，未重新取得十四份分卷字节。MIT 核验页公开提供源码镜像，并说明原源码与镜像相同，镜像后续另加比较器文件。本轮固定恢复社区仓库的首提交
\begin{quote}\small
\path{40e56bab026552e089d7afcfeff08d2defb9c005}
\end{quote}
的完整源码树，不把 GitHub 重新封装的 gzip 称为作者原始 gzip。先前保留的十四分卷 SHA256 清单同时附上，并标记本轮没有重新逐分卷校验。

实际恢复的项目包含 $175\,187$ 个 Lean 源文件，正文共 $1\,030\,531\,854$ 字节、$17\,374\,215$ 个物理行，另有三个项目配置文件。本轮将先前保留的 $88$ 个原源码及配置文件逐字节对照，全部一致。这一结果检查的是源码文本和来源，不能代替 Lean 内核对证明的检查。

\section{单文件汇编保留了什么}
交付文件 \path{RH_AllSources.lean} 按项目内部导入关系的拓扑顺序收录全部模块，在每个模块前后加入可识别的注释边界。各模块原始正文的字节保持不变，没有删除证明体，也没有用 \texttt{sorry} 或新增公理替代证明。

配套的 \path{RH_AllSources.index.jsonl} 记录模块名、原路径、正文的字节范围、行数和 SHA256；\path{RH_AllSources.dependencies.jsonl} 记录直接导入边。恢复工具读取汇编文件自身包含的模块边界、长度和哈希，抽取原文并放回原项目路径；恢复时不要求另有外部索引。因此保留了逐模块查阅和恢复编译工程的能力。直接导入图表示源文件依赖；它不等同于内核导出的最小定理依赖图。

\textbf{编译性质必须分清。}这个文件是无损源码汇编，尚未被改造为可独立编译的单一 Lean 模块。Lean 规定导入位于文件头；模块名还参与私有声明和作用域处理。把许多文件的正文拼接起来，会在正文中再次出现导入，并改变原来彼此隔离的局部环境。仅仅删除导入、给所有文件套一个命名空间，也不能保证保留原证明的含义。为此，本轮保留原文并提供可逆恢复，没有把文本合并的成功写成新单体证明的编译成功。

原项目的 Mathlib 与 Lean 工具链仍按 \path{lean-toolchain}、\path{lakefile.toml}、\path{lake-manifest.json} 固定。合并十四分卷所含源码，不会自动把 Lean 编译器、Mathlib 及全部外部依赖也变成同一个文件。

\clearpage
\section{三种证明文件各自说明什么}
\begin{longtable}{@{}p{.24\textwidth}p{.69\textwidth}@{}}
\toprule
文件种类 & 可以回答的问题\\
\midrule
\endhead
完整源码汇编 & 原作者写了哪些定义、定理和证明？某个模块的原文在哪里？能否无损恢复？\\
数学说明与关键源码 & 从目标声明到算术窗口、解析排零点和函数方程，各步骤如何相接？前提是什么？对应哪段代码？\\
内核证明项导出 & 展开后的声明与证明项是什么？独立检查器能否逐项接受？\\
\bottomrule
\end{longtable}

本轮实际导出的《证明过程\_数学说明与关键源码.txt》和《证明过程\_结构化索引.json》按十三个阶段组织说明，保留三十三个关键模块的原文、逐行编号及哈希。这个过程文件是源码级说明；没有把静态文本抽取冒充 Lean 执行时的逐 tactic 状态记录。

MIT 页面另公开提供完整内核导出 \path{solution.ndjson.zst}：压缩文件为 $7\,699\,933\,588$ 字节，解压为 $61\,570\,595\,788$ 字节。它与约一 GB 的 Lean 源码是不同层次的对象。本轮没有下载并复验该大文件；附带的来源记录给出其正式链接、哈希和公开核验入口，以便在具备相应资源的环境继续。不能把本文的过程说明当作该内核导出的替代品。

\section{如何沿证明过程继续阅读}
先从过程文件的 P01 核对最终命题，再沿 P03--P08 阅读窗口的算术来源，随后由 P09--P12 检查完整解析推导，最后在 P13 核对函数方程与特殊点。若只想理解正文的数学结构，可先读第\ref{chap:overview}章的记号表及第\ref{chap:window}章的上下界矛盾，再回到第\ref{chap:arithmetic}章逐项核对输入。

本次比上一版增加了上同调长度的有限部分和解释、宿主请求的量词、平方窗口缩到 $3/2$ 次方窗口、$401\to400$ 的截断余量、有限特征正交与 $4o$ 计数，以及统一阈值的拼接。尚需继续转写的是这些接口更深层的自守、Galois、几何和模性构造；源码完整收录不意味着这些构造已经全部成为人类文字证明。
